\chapter{Literature Review}
\label{ch:literature}

\section{Deep Learning and Vision Transformers}

Deep learning has revolutionized almost every aspect of artificial intelligence, with groundbreaking success in practically all of computer vision and natural language processing. Goodfellow et al.~\cite{goodfellow2016deep} give a maths heavy account of neural networks, including optimization methods, regularization techniques, and architectural ‘hierarchies’ which give deep networks their representational power. For image tasks, deep residual networks were introduced to combat the vanishing gradient problem using identity shortcut connections \cite{he2016deep}, enabling extremely deep networks to be trained and achieve record results on huge benchmarks such as ImageNet. 

The Vision Transformer (ViT) shifted attention from CNNs to attention based models by demonstrating that a pure transformer architecture over a sequence of image patches can be competitive with convolutional architectures if sufficiently large \cite{dosovitskiy2020image}. Unlike CNNs, which implicitly encode inductive biases of spatial locality and translation equivariance into their architectures, ViT models are required to learn from scratch to attend to spatial aspects of an image through self-attention, providing them with the capacity to capture long-range dependencies of particular use in forensic tasks for which artifact signatures may lie in non-local parts of the image.

A systematic layer-wise investigation of CLIP-based Vision Transformers for AI-generated image detection by Park et al.~\cite{park2025rethinking} shows that intermediate layers of ViT models rather than the final classification layer provide the most generalizable forensic discrimination features. They propose a Mixture of Layers for Detection (MoLD) mechanism that dynamically combines features across several ViT layers through a learned gating mechanism, yielding better generalization to previously unseen generative models. This observation supports the theory that intermediate feature representations acquired during vision-language pre-training encode forensically meaningful characteristics that are lost in last-layer compression.

\section{Vision-Language Models and SigLIP 2}

The intersection of visual and linguistic representation learning has led to a family of capable foundation models with exceptional cross-modal understanding. CLIP \cite{radford2021learning} established the format of contrastive vision-language pre-training, aligning image and text representations through a contrastive objective that pulls similar image-text pairs to share similar embeddings. Trained on 400 million image-text pairs, the resulting model demonstrated strong zero-shot transfer to image classification tasks.

Liu et al.~\cite{liu2024fatformer} proposed FatFormer, a forgery-aware adaptive transformer that explicitly utilises the vision-language spaces of pre-trained CLIP for generalised synthetic image detection. FatFormer combines local forgery traces in both image and frequency domains using a forgery-aware adapter module, and uses language-guided alignment to regularise the learned features. The approach scored 98\% average accuracy on unseen GAN architectures and 95\% on unseen diffusion models, demonstrating the value of vision-language pre-training as a foundation for forensic detection.

SigLIP 2, proposed by Tschannen et al.~\cite{tschannen2025siglip2}, is a significant improvement to CLIP-based architectures. Rather than using a contrastive loss with softmax normalisation that requires large batch sizes for effective training, SigLIP 2 adopts a pairwise sigmoid loss where each image-text pair is trained independently. The model was trained on the WebLI dataset comprising 10 billion image-text pairs and supports distortion-free native image resolutions through a flexible resolution handling system named NaFlex. These properties combine to form a vision encoder with highly dense and semantically rich feature representations---characteristics that make it well suited to the fine-grained discrimination required in image forensics.

\section{AI-Generated Image Detection: State of the Art}

The challenge of distinguishing AI-generated from real images has evolved significantly since early GAN detection work. Bonettini and Bestagini~\cite{bonettini2020benford} capitalised on the statistical characteristics of Benford's Law in the DCT coefficient distributions of GAN-generated images, observing that the natural distribution of significant digits of real camera images is not replicated by GAN synthesis. This frequency-domain perspective provided a significant foundation for the forensic feature engineering approach used in the present work.

Diffusion model-based generators rendered the detection landscape considerably more difficult. Li et al.~\cite{li2025aigibench} proposed AIGIBench, a large benchmark comprising 23 diverse subsets of AI-generated images from a variety of generative models, compared against real images from FFHQ, CelebA-HQ, and Open Images V7. Their evaluation of 11 state-of-the-art detectors demonstrated that all current methods exhibit steep accuracy deterioration under real-world conditions including JPEG compression, image resizing, and colour jitter, with some detectors declining to near-random performance under modest perturbation.

Ren et al.~\cite{ren2026benchmark} extended these findings with the first large-scale evaluation of 16 detection approaches in a zero-shot setting across 12 diverse datasets, including GenImage, WildFake, OpenFake, and SynthBuster. Their results revealed enormous performance variation across datasets, with the model performing best on one benchmark potentially performing worst on another, and proposed that training data compatibility is a more important performance predictor than algorithmic design choices. These results encourage the creation of more eclectic training datasets and more abstract feature representations.

Perotti et al.~\cite{perotti2025nodedetector} demonstrated that available open-source deepfake detectors are not capable of generalising across different generative models. Through extensive experimentation, they demonstrated that ensemble methods with latent feature fusion are reliable for improving single binary classifiers. However, the authors discussed the fact that ensemble advantages are acquired at high computational cost and thus efficient single model alternatives are important.

Karageorgiou et al.~\cite{karageorgiou2024spai} created SPAI (Spectral AI generated Image detection), an approach based on masked spectral learning to identify AI images as out-of-distribution samples by determining their similarity to a spectral reconstruction distribution that has been trained on real images. Evaluated on generative models not seen during training, SPAI improved accuracy over state-of-the-art methods by 5.5\%, demonstrating the strength of spectral domain analysis for generalized detection.

Zhong et al.~\cite{zhong2024patchcraft} proposed PatchCraft, a preprocessing scheme that exploits differences in inter-pixel correlation between rich- and poor-texture regions of AI-generated images. PatchCraft disrupts global semantic coherence by shuffling texture patches, forcing a detector to attend to local texture statistics, achieving an average zero-shot accuracy of 89.85\% across six generative models.

\section{Hybrid Architectures for Efficient Detection}

Beyond pure approaches of deep learning, a range of studies have been conducted on hybrid solutions which mix the power of deep learning with the interpretability and hand-crafted features. Lokner Ladi’c et al.~\cite{loknerladic2023cnn} proposed a model which jointly uses a CNN along with a wavelet entropy technique, achieving 97.32\% accuracy whilst remaining computationally acceptable due to being simpler. Their explainability analysis showed that in this architecture the model focuses primarily on certain complex textures of the kind found in backgrounds, which is very similar to the LBP and gradient co-occurrence forensic signals used in our work.

Gupta and McMahan~\cite{gupta2025fft} showed that preprocessing the input images fed to a CNN based architecture with FFT leads to +3.27\% improved accuracy and 10.06\% reduced loss compared to models that do not use frequency-domain features, justifying our use of the FFT based power slope forehead signal in the present architecture.

Rani et al.~\cite{rani2025lightweight} present a model using VisionTransformers and Linformer structure developed to detect deepfake videos with an accuracy of 98.9\% whilst being more computationally efficient than the standard ViT structures. Harshit et al.~\cite{harshit2025deepfake} have in fact presented more recently a model also using EfficientNet-B0, along with a GRU, used to detect deepfake videos.

\section{Traditional Forensic Signals in the Deep Learning Era}

A growing area of research is focused on the exploration of different approaches to integrate the available traditional data features into deep learning models. It is a known fact that modern deep learning architectures are plagued with a number of limitations. Tan et al.~\cite{tan2023gradients} showed that gradient maps are a universal representation for detecting GAN-generated images, improving performance by 11.4% compared to previous methods, where all images are transformed into the gradient domain through a fixed differential operator. Thus, the transformation domain is directly related to the Gradient Co-occurrence feature used in this paper.

Significant research has been conducted in photo forensics, studying camera sensor noise as a forensic signal. Photographs, during Bayer demosaicing, have inherent correlations between color channel noise residuals, which are absent in AI-generated images, since color values in all three channels are synthesized simultaneously. Hence, a discriminative signal, robust against post-processing operations affecting the physical content of images, has been obtained.

The Local Binary Pattern (LBP) was originally proposed by Ojala et al.~\cite{ojala2002lbp} for texture classification, which has been extended significantly for image authenticity verification. The use of Shannon entropy values for LBP has been shown to be effective, since a scalar value representing statistical distribution of micro-textures in images has been obtained. Photographs, being an inhomogeneous mixture of edges, gradients, sensor noise, and physical surface textures, tend to have LBP entropy values closer to the theoretical limit, while AI-generated images tend to have LBP values that systematically deviate from those obtained from real photographs, since AI-generated images, which aim to be coherent at a macroscopic level, tend to be poor in physical realism at a micro-level. Hence, LBP has been proposed as a discriminative feature, distinguishing between AI-generated images and real photographs, with an area under the curve (AUC) value of 0.79 during the feature selection phase of this research, making LBP the best discriminative feature among traditionally used features.

\section{Literature Gap and Research Positioning}

A comprehensive review of the existing literature reveals a critical gap at the intersection of three research directions: the use of VLM-pretrained vision encoders with a sigmoid pairwise loss (SigLIP 2), the explicit combination of deep embeddings with traditional forensic signals in a unified hybrid architecture, and the rigorous evaluation of such an architecture against hard-negative images and commercial detection baselines under realistic deployment constraints.

While FatFormer \cite{liu2024fatformer} demonstrated the value of CLIP-based vision-language pre-training for detection, it does not explore the SigLIP 2 encoder, does not incorporate traditional forensic signals, and does not report inference latency under web-deployment conditions. While Gupta and McMahan~\cite{gupta2025fft} and Tan et al.~\cite{tan2023gradients} demonstrated the value of frequency-domain and gradient features respectively, neither work combined these signals with a VLM-quality vision backbone. The benchmark studies of Li et al.~\cite{li2025aigibench} and Ren et al.~\cite{ren2026benchmark} establish the severity of the generalisation problem but offer no architectural solutions. The present dissertation directly addresses this gap by proposing, implementing, and rigorously evaluating a hybrid architecture that, to the best of the author's knowledge, is the first to combine a SigLIP 2 vision encoder with an explicit traditional forensic feature branch for AI image detection.
